\documentclass[final]{beamer}

\usepackage[size=a0,orientation=landscape,scale=1.03]{beamerposter}
\usepackage[T1]{fontenc}
\usepackage[utf8]{inputenc}
\usepackage{lmodern}
\usepackage{booktabs}
\usepackage{array}
\usepackage{tikz}
\usepackage{ragged2e}

\definecolor{Ink}{HTML}{172019}
\definecolor{Muted}{HTML}{5B645D}
\definecolor{Paper}{HTML}{F7F0DF}
\definecolor{Cream}{HTML}{FFFAF0}
\definecolor{Green}{HTML}{0F7F64}
\definecolor{Clay}{HTML}{C8643D}
\definecolor{Blue}{HTML}{315F85}
\definecolor{Gold}{HTML}{D4A72C}
\definecolor{Charcoal}{HTML}{18221D}
\definecolor{SoftGreen}{HTML}{D8EFE3}
\definecolor{SoftBlue}{HTML}{D7E5EE}
\definecolor{SoftGold}{HTML}{F5E6B5}

\setbeamercolor{background canvas}{bg=Paper}
\setbeamercolor{normal text}{fg=Ink}
\setbeamercolor{block title}{fg=Cream,bg=Charcoal}
\setbeamercolor{block body}{fg=Ink,bg=Cream}
\setbeamerfont{block title}{size=\large,series=\bfseries}
\setbeamerfont{block body}{size=\small}
\setbeamertemplate{navigation symbols}{}
\setbeamertemplate{headline}{}
\setbeamertemplate{footline}{}

\newcommand{\postertag}[1]{%
  \tikz[baseline=(X.base)]\node[
    draw=Ink!18,
    fill=Cream,
    rounded corners=7pt,
    inner xsep=7pt,
    inner ysep=3pt,
    font=\bfseries\footnotesize
  ] (X) {#1};%
}

\newcommand{\hbarval}[5]{%
  \begin{tikzpicture}[baseline=-0.5ex, x=0.095cm, y=0.55cm]
    \node[anchor=east, font=\bfseries\scriptsize] at (0,0) {#1};
    \fill[Ink!10, rounded corners=2pt] (1,-0.22) rectangle (101,0.22);
    \fill[#3, rounded corners=2pt] (1,-0.22) rectangle ({1+#2},0.22);
    \node[anchor=west, font=\bfseries\scriptsize] at (104,0) {#4};
    \node[anchor=west, font=\scriptsize, text=Muted] at (118,0) {#5};
  \end{tikzpicture}%
}

\newcommand{\progressbar}[5]{%
  \begin{tikzpicture}[baseline=-0.5ex, x=0.11cm, y=0.55cm]
    \node[anchor=east, font=\bfseries\scriptsize] at (0,0) {#1};
    \fill[Ink!10, rounded corners=2pt] (1,-0.22) rectangle (101,0.22);
    \fill[#3, rounded corners=2pt] (1,-0.22) rectangle ({1+100*#2/148},0.22);
    \node[anchor=west, font=\bfseries\scriptsize] at (104,0) {#2/148};
    \node[anchor=west, font=\scriptsize, text=Muted] at (122,0) {#4 capped; #5};
  \end{tikzpicture}%
}

\newcommand{\darkbody}{\setbeamercolor{block body}{fg=Cream,bg=Charcoal}}
\newcommand{\lightbody}{\setbeamercolor{block body}{fg=Ink,bg=Cream}}

\title{TA-MPQ: Structured Mixed Precision at an Exact INT4 Budget}
\author{Snapshot: 2026-04-22}
\institute{Qwen3.5-9B \quad | \quad Current mainline: deterministic 2/4/8-bit frontier search}

\begin{document}
\begin{frame}[t]

\begin{tikzpicture}[remember picture,overlay]
  \fill[Gold!20] (current page.north west) circle [radius=18cm];
  \fill[Blue!14] ([xshift=-8cm,yshift=-2cm]current page.north east) circle [radius=22cm];
  \draw[Ink!8, line width=2pt, rounded corners=18pt]
    ([xshift=1.0cm,yshift=-1.0cm]current page.north west)
    rectangle
    ([xshift=-1.0cm,yshift=1.0cm]current page.south east);
\end{tikzpicture}

\vspace{-0.2cm}
{\fontsize{58}{62}\selectfont\bfseries\color{Charcoal}\inserttitle\par}
\vspace{0.25cm}
{\Large\color{Muted}
Replacing evolutionary crossover with a task-sensitivity frontier: promote high-value groups to 8-bit,
demote lower-value groups to 2-bit, and keep the full model near the raw uniform-INT4 footprint.\par}
\vspace{0.25cm}
\postertag{Qwen3.5-9B}\hspace{0.18cm}
\postertag{MMLU-coding selected}\hspace{0.18cm}
\postertag{EvalPlus validated}\hspace{0.18cm}
\postertag{BigCodeBench generation running}\hspace{0.18cm}
\postertag{Budget = raw weight footprint}

\vspace{0.55cm}

\begin{columns}[t,totalwidth=\textwidth]
  \begin{column}{0.245\textwidth}
    \begin{block}{Motivation and route change}
      \justifying
      The original TA-MPQ implementation used evolutionary search over mixed-precision policies.
      That approach was difficult to control because crossover randomly combined policies while
      the budget mechanism tried to repair them back to a strict footprint target. These two forces
      were misaligned: crossover explores, but exact-budget quantization needs coordinated
      promotions and demotions.

      \vspace{0.28cm}
      The current route is deterministic. We start from the uniform INT4 anchor and build an ordered
      frontier. Moving along the frontier adds INT8 to valuable groups and offsets that extra cost by
      assigning INT2 to lower-value groups. This makes the policy easier to debug, reproduce, and
      explain.

      \vspace{0.28cm}
      \colorbox{SoftGreen}{\parbox{0.95\linewidth}{
      \textbf{Poster claim.} The new route is not yet a universal win. It is a controllable
      INT4-sized policy that recovers part of the INT4 quality loss on generative code tasks.
      }}
    \end{block}

    \begin{block}{Pipeline}
      \begin{tabular}{@{}p{0.23\linewidth}p{0.71\linewidth}@{}}
      \textbf{Profile} & Collect task sensitivity on MMLU-coding development prompts.\\[0.13cm]
      \textbf{Rank} & Estimate which groups deserve 8-bit and which groups can tolerate 2-bit.\\[0.13cm]
      \textbf{Frontier} & Build a greedy path under exact raw INT4 footprint.\\[0.13cm]
      \textbf{Refine} & Sample coarse regions, then refine near the best local neighborhood.\\[0.13cm]
      \textbf{Validate} & Quantize the selected policy into a real artifact and evaluate held-out tasks.\\
      \end{tabular}

      \vspace{0.28cm}
      Current selected candidate: \texttt{gpf\_refine\_00\_i0112}.\\
      Policy hash: \texttt{7caae6e...2043e4}.\\
      Method: \texttt{greedy\_path\_frontier}.\\
      Search space: \texttt{\{2,4,8\}}.
    \end{block}
  \end{column}

  \begin{column}{0.245\textwidth}
    \darkbody
    \begin{block}{Current mixed policy composition}
      \justifying
      The selected mixed model changes 169 of 249 groups relative to uniform INT4 while matching the
      INT4 raw weight budget. The policy is not simply ``many groups in INT8.'' It keeps most parameter
      mass at 4-bit, upgrades a smaller high-value mass to 8-bit, and pays for those upgrades using
      2-bit assignments on lower-value groups.

      \vspace{0.35cm}
      \begin{tikzpicture}
        \fill[Clay, rounded corners=5pt] (0,0) rectangle (3.12,0.48);
        \fill[Gold] (3.12,0) rectangle (10.44,0.48);
        \fill[Blue, rounded corners=5pt] (10.44,0) rectangle (12,0.48);
        \node[anchor=west, text=Cream, font=\bfseries\scriptsize] at (0,0.85) {Param mass: 26.00\% 2-bit};
        \node[anchor=west, text=Cream, font=\bfseries\scriptsize] at (4.0,0.85) {61.02\% 4-bit};
        \node[anchor=west, text=Cream, font=\bfseries\scriptsize] at (8.5,0.85) {12.98\% 8-bit};
      \end{tikzpicture}

      \vspace{0.35cm}
      \begin{tabular}{lrr}
        \toprule
        Bit width & Groups & Parameter mass\\
        \midrule
        2-bit & 57 & 26.00\%\\
        4-bit & 80 & 61.02\%\\
        8-bit & 112 & 12.98\%\\
        \bottomrule
      \end{tabular}

      \vspace{0.25cm}
      \begin{tabular}{lr}
        Total groups & 249\\
        Changed vs uniform INT4 & 169\\
        Estimated average bit-width & 3.9989\\
        Estimated raw footprint & 3.694 GB\\
        Budget slack & 0.026\%\\
      \end{tabular}
    \end{block}
    \lightbody

    \begin{block}{MMLU-coding chart}
      \small Simple-evals style prompt, max\_new\_tokens=4096, single-run last100.

      \vspace{0.25cm}
      \hbarval{BF16}{95}{Green}{95\%}{0 capped}\\[0.18cm]
      \hbarval{INT8}{94}{Blue}{94\%}{2 capped}\\[0.18cm]
      \hbarval{INT4}{91}{Clay}{91\%}{5 capped}\\[0.18cm]
      \hbarval{Mixed}{91}{Gold}{91\%}{3 capped}

      \vspace{0.25cm}
      \begin{tabular}{lrrr}
        \toprule
        Model & Correct & Accuracy & Avg latency\\
        \midrule
        BF16 & 95/100 & 95\% & 9.14s\\
        INT8 & 94/100 & 94\% & 46.21s\\
        INT4 & 91/100 & 91\% & 66.88s\\
        Mixed & 91/100 & 91\% & 12.29s\\
        \bottomrule
      \end{tabular}
    \end{block}
  \end{column}

  \begin{column}{0.245\textwidth}
    \begin{block}{HumanEval / EvalPlus chart}
      \small Official EvalPlus evaluator, greedy generation, max\_new\_tokens=768.

      \vspace{0.25cm}
      \begin{tikzpicture}[x=0.11cm,y=0.065cm]
        \draw[Ink!18] (0,0) -- (100,0);
        \foreach \x in {0,20,40,60,80,100} {
          \draw[Ink!12] (\x,0) -- (\x,62);
          \node[font=\tiny,text=Muted] at (\x,-4) {\x\%};
        }
        \node[anchor=east,font=\scriptsize\bfseries] at (-3,55) {BF16};
        \fill[Green] (0,52) rectangle (87.8,58);
        \fill[Green!55] (0,45) rectangle (82.3,51);
        \node[anchor=west,font=\tiny] at (89,55) {87.8};
        \node[anchor=west,font=\tiny] at (83.5,48) {82.3};

        \node[anchor=east,font=\scriptsize\bfseries] at (-3,38) {INT8};
        \fill[Blue] (0,35) rectangle (89.0,41);
        \fill[Blue!55] (0,28) rectangle (82.3,34);
        \node[anchor=west,font=\tiny] at (90.2,38) {89.0};
        \node[anchor=west,font=\tiny] at (83.5,31) {82.3};

        \node[anchor=east,font=\scriptsize\bfseries] at (-3,21) {INT4};
        \fill[Clay] (0,18) rectangle (81.1,24);
        \fill[Clay!55] (0,11) rectangle (76.8,17);
        \node[anchor=west,font=\tiny] at (82.3,21) {81.1};
        \node[anchor=west,font=\tiny] at (78,14) {76.8};

        \node[anchor=east,font=\scriptsize\bfseries] at (-3,4) {Mixed};
        \fill[Gold] (0,1) rectangle (84.8,7);
        \fill[Gold!58] (0,-6) rectangle (79.9,0);
        \node[anchor=west,font=\tiny] at (86,4) {84.8};
        \node[anchor=west,font=\tiny] at (81.1,-3) {79.9};

        \fill[Ink] (4,66) rectangle (9,69);
        \node[anchor=west,font=\tiny] at (11,67.5) {HumanEval base};
        \fill[Ink!45] (40,66) rectangle (45,69);
        \node[anchor=west,font=\tiny] at (47,67.5) {HumanEval+};
      \end{tikzpicture}

      \vspace{0.2cm}
      \begin{tabular}{lrr}
        \toprule
        Model & Base pass@1 & Plus pass@1\\
        \midrule
        BF16 & 87.8\% & 82.3\%\\
        INT8 & 89.0\% & 82.3\%\\
        INT4 & 81.1\% & 76.8\%\\
        \textbf{Mixed} & \textbf{84.8\%} & \textbf{79.9\%}\\
        \bottomrule
      \end{tabular}

      \vspace{0.22cm}
      \colorbox{SoftGreen}{\parbox{0.95\linewidth}{
      \textbf{Transfer signal:} mixed recovers +3.7 points over INT4 on HumanEval base
      and +3.0 points on HumanEval+, while staying at the INT4-sized raw footprint.
      }}
    \end{block}

    \begin{block}{BigCodeBench progress chart}
      \small Official BigCodeBench instruct prompt. Generation-only first; execution grading pending.

      \vspace{0.25cm}
      \progressbar{BF16}{148}{Green}{2}{complete}\\[0.18cm]
      \progressbar{INT8}{20}{Blue}{0}{H100 running}\\[0.18cm]
      \progressbar{INT4}{20}{Clay}{0}{H100 running}\\[0.18cm]
      \progressbar{Mixed}{30}{Gold}{1}{H100 running}

      \vspace{0.25cm}
      We intentionally do not report BigCodeBench accuracy yet. The previous local evaluator timed out
      because the execution environment was incomplete. The current runner writes generation artifacts
      before evaluation, so evaluator failures will not erase samples.
    \end{block}
  \end{column}

  \begin{column}{0.245\textwidth}
    \darkbody
    \begin{block}{Current interpretation}
      \begin{itemize}
        \item The current story is not ``mixed beats every baseline.''
        \item The current story is that structured exact-budget mixed precision is a more defensible mainline than evolutionary crossover.
        \item INT8 remains a very strong quantized upper baseline.
        \item Uniform INT4 is size-efficient but loses quality on generative coding tasks.
        \item Mixed recovers part of that INT4 degradation on HumanEval/EvalPlus without increasing the raw weight budget.
        \item MMLU-coding does not yet separate mixed from INT4; it is not the strongest success case.
        \item BigCodeBench is the next decisive benchmark because it is generative, dependency-heavy, and closer to realistic code synthesis.
      \end{itemize}
    \end{block}
    \lightbody

    \begin{block}{Limitations}
      \begin{itemize}
        \item MMLU-coding last100 is a small single-run slice.
        \item HumanEval has only 164 tasks, so effect sizes should be interpreted carefully.
        \item BigCodeBench execution scoring is still pending.
        \item The budget here is raw weight footprint, not peak VRAM, kernel efficiency, or end-to-end serving cost.
        \item The mixed policy was selected using MMLU-coding; stronger cross-task validation is needed before making a broad task-aware claim.
      \end{itemize}
    \end{block}

    \begin{block}{Next steps}
      \begin{enumerate}
        \item Finish BigCodeBench generation for INT8, INT4, and mixed.
        \item Run official BigCodeBench execution grading in a stable environment.
        \item Report peak VRAM and tokens/sec separately from raw weight footprint.
        \item Use BigCodeBench results to decide whether MMLU-coding sensitivity is sufficient or needs task-specific refinement.
        \item If BigCodeBench shows mixed gains, promote this structured exact-budget route as the main TA-MPQ story.
        \item If BigCodeBench does not show mixed gains, improve the sensitivity profiler and candidate generator before claiming task-aware benefit.
      \end{enumerate}
    \end{block}
  \end{column}
\end{columns}

\vfill
\begin{center}
{\color{Muted}\large
Recommended conclusion: Structured exact-budget mixed precision is a real, reproducible INT4-sized quantized model that recovers part of the INT4 degradation on generative code evaluation. The next decisive result is BigCodeBench execution accuracy.}
\end{center}

\end{frame}
\end{document}
